\chapter{Implementation}
\label{ch:implementation}

This chapter describes how the architecture of Chapter~\ref{ch:architecture} is
realised in code. It moves from the development stack and repository layout, down
through the on-chain contracts, and back up through the client-side cryptography,
threshold key custody, permanent storage, anonymous payments, and zero-knowledge
provenance, finishing with the way these pieces are orchestrated in the browser
and wired together at deployment time. Throughout, the emphasis is on the design
decisions that give \zkw its security properties --- where a secret lives, who can
read it, and what an adversary observes --- rather than on incidental UI glue.

The complete source is organised as a monorepo; the listings in this chapter are
faithful extracts, occasionally abridged with comments to keep the discussion
focused. File paths are given relative to the repository root so that each listing
can be located in the accompanying code.

\section{Implementation Overview}
\label{sec:impl-overview}

\zkw is built on Scaffold-ETH~2 (SE-2), a TypeScript starter kit that pairs a
Hardhat smart-contract workspace with a Next.js front end and a set of typed
React hooks for contract interaction. The choice is deliberate: SE-2 provides a
reproducible local chain, automatic generation of contract ABIs and typed
bindings, and a component library for common Web3 primitives, which lets the
project concentrate engineering effort on the parts that are genuinely novel --- the
cryptographic pipeline and the trust model --- rather than on boilerplate.

\begin{table}[h]
\centering
\caption{Implementation stack by layer.}
\label{tab:impl-stack}
\begin{tabular}{p{3.6cm} p{4.0cm} p{6.0cm}}
\toprule
\textbf{Layer} & \textbf{Technology} & \textbf{Role in \zkwns} \\
\midrule
Contracts      & Solidity 0.8.x, Hardhat, OpenZeppelin & Dead man's switch, registry, marketplace, on-chain verification \\
Client crypto  & Web Crypto API (browser) & AES-256-GCM encryption of payloads, all in-browser \\
Key custody    & Lit Protocol SDK v8 (Naga) & Threshold sealing/unsealing of the AES key under on-chain conditions \\
Provenance     & Reclaim Protocol (zkTLS) & Anonymous proof of source credentials \\
Payments       & \code{@noble/secp256k1} & ERC-5564 stealth-address derivation and scanning \\
Storage        & Arweave via Turbo SDK & Permanent, censorship-resistant payload storage \\
Front end      & Next.js (App Router), wagmi, viem, DaisyUI & Wallet connection, vault wizard, release flow \\
\bottomrule
\end{tabular}
\end{table}

The repository is split into two packages under \code{packages/}: \code{hardhat}
holds the Solidity sources, deployment scripts, and tests; \code{nextjs} holds the
front end, including a \code{services/zk-whistle} directory that contains the
security-critical, framework-agnostic logic (encryption, Lit custody, stealth
addresses, storage, and Reclaim). Keeping this logic in plain TypeScript modules,
separate from React components, makes it independently testable and auditable, and
keeps the cryptographic core free of UI concerns.

A guiding implementation rule, enforced by construction rather than convention, is
that \emph{no plaintext and no key material ever leaves the browser in the clear}.
Every network egress point --- the storage route, the chain, the Lit network --- only
ever sees ciphertext or values derived from it. The sections that follow show how
each component upholds that rule.

\section{The Smart-Contract Layer}
\label{sec:impl-contracts}

Three contracts make up the on-chain logic. \code{DeadMansSwitch} is the trigger
oracle; \code{WhistleblowerRegistry} anchors anonymous reputation; and
\code{Marketplace} settles payments to stealth addresses. A small interface,
\code{IReclaim}, lets the registry delegate cryptographic proof checking to a
deployed Reclaim verifier. All three contracts deliberately store \emph{no}
sensitive data: they hold timestamps, hashes, references, and payment state, never
plaintext or keys.

\subsection{The Dead Man's Switch}
\label{sec:impl-dms}

The dead man's switch is intentionally small. Its public surface is a single
mapping from an owner address to a \code{Switch} struct, a handful of
state-changing functions (\code{createSwitch}, \code{checkIn},
\code{deactivateSwitch}, \code{updateSwitchMetadata}), and the read function that
the entire release mechanism hinges on: \isdeceased. Listing~\ref{lst:dms-core}
shows the heart of the contract.

\begin{lstlisting}[language=Solidity,caption={The switch struct and the trigger predicate (\code{DeadMansSwitch.sol}).},label={lst:dms-core}]
struct Switch {
    uint256 lastHeartbeat;
    uint256 heartbeatInterval; // seconds between required check-ins
    string  arweaveTxId;       // Arweave tx id of encrypted payload
    string  litAccessControlId;// Lit access-control reference
    address recipient;         // address(0) = public release
    bool    isActive;
}

mapping(address => Switch) public switches;

function createSwitch(
    uint256 _heartbeatInterval,
    string calldata _arweaveTxId,
    string calldata _litAccessControlId,
    address _recipient
) external {
    require(!switches[msg.sender].isActive, "Switch already active");
    require(_heartbeatInterval > 0, "Interval must be > 0");
    require(bytes(_arweaveTxId).length > 0, "Arweave TX ID required");

    switches[msg.sender] = Switch({
        lastHeartbeat: block.timestamp,
        heartbeatInterval: _heartbeatInterval,
        arweaveTxId: _arweaveTxId,
        litAccessControlId: _litAccessControlId,
        recipient: _recipient,
        isActive: true
    });
    // ... enumeration bookkeeping and events ...
}

function checkIn() external onlyActiveSwitchOwner {
    switches[msg.sender].lastHeartbeat = block.timestamp;
    emit HeartbeatUpdated(msg.sender, block.timestamp);
}

function isDeceased(address _user) external view returns (bool) {
    Switch storage s = switches[_user];
    if (!s.isActive) return false;
    return block.timestamp > s.lastHeartbeat + s.heartbeatInterval;
}
\end{lstlisting}

Three properties of this design matter for the rest of the system. First,
\isdeceased is a pure function of public on-chain state --- a timestamp comparison ---
so any party, including the Lit network's nodes, can evaluate it independently and
arrive at the same answer; the trigger cannot be faked or suppressed by a single
actor. Second, deactivation makes a switch return \code{false} even after the
interval elapses, so a source who is safe can permanently disarm their vault.
Third, the contract stores only an \code{arweaveTxId} and a
\code{litAccessControlId}: pointers to where the ciphertext and the sealed key
live, never the secrets themselves. This is what makes the on-chain footprint of a
vault information-free with respect to its contents.

The companion view \code{timeUntilTrigger} returns the seconds remaining before a
switch fires, which the front end uses to render a countdown and to nudge the
source to check in. It returns zero for inactive or already-triggered switches, so
the UI never has to special-case those states.

\subsection{The Whistleblower Registry and Tiered Verification}
\label{sec:impl-registry}

The registry stores a profile per address: a list of proof hashes, a verified
flag, a registration timestamp, and an ERC-5564 stealth meta-address. The subtlety
is in what ``verified'' is allowed to mean. A na\"ive implementation that lets a
user self-assert verification would make the badge worthless. \zkw therefore
offers three verification tiers, ordered from strongest to weakest, and the
deployment chooses which are active.

The strongest tier is fully trustless on-chain verification. The caller submits a
complete Reclaim proof; the registry forwards it to a deployed Reclaim verifier,
which checks the witness signatures on-chain. Crucially, verification accrues to
the proof's \emph{owner} --- the address the proof was issued to --- not to the caller,
which prevents a valid proof from being replayed to verify a different account and
allows the call to be relayed by a gasless sponsor on the owner's behalf
(Listing~\ref{lst:registry-verify}).

\begin{lstlisting}[language=Solidity,caption={Trustless on-chain proof verification (\code{WhistleblowerRegistry.sol}).},label={lst:registry-verify}]
function submitVerifiedProof(ReclaimTypes.Proof calldata proof) external {
    require(address(reclaimVerifier) != address(0), "Verifier not configured");

    bytes32 identifier = proof.signedClaim.claim.identifier;
    address subject    = proof.signedClaim.claim.owner;
    require(identifier != bytes32(0), "Invalid proof identifier");
    require(subject != address(0),    "Invalid proof owner");
    require(!proofHashExists[identifier], "Proof already submitted");

    // Reverts if witness signatures are invalid or below threshold.
    reclaimVerifier.verifyProof(proof);

    _ensureRegistered(subject);
    profiles[subject].proofHashes.push(identifier);
    proofHashExists[identifier] = true;

    if (!profiles[subject].isVerified) {
        profiles[subject].isVerified = true;
        emit VerificationUpdated(subject, true);
    }
    emit ProofVerifiedOnChain(subject, identifier);
}
\end{lstlisting}

The middle tier, \code{attestVerification}, lets an authorised verifier (the
contract owner, standing in for an off-chain validator or a Lit Action) set a
user's status after checking a proof off-chain. It exists for chains or providers
where a Reclaim verifier is not deployed. The weakest tier,
\code{autoVerifyOnSubmit}, trust-marks a user on their first proof hash; it is a
development convenience only and is enabled exclusively on the local chain (see
Section~\ref{sec:impl-deploy}). Encoding the trust assumption directly in the
contract, and making the production path cryptographic, keeps the meaning of the
``Verified Source'' badge honest.

\subsection{The Marketplace and Anonymous Settlement}
\label{sec:impl-marketplace}

The marketplace lets sources monetise information without ever revealing who
receives the money. Listings carry a description hash, an Arweave reference, a
minimum bid, and --- importantly --- a verified flag that is \emph{read from the
registry at creation time} rather than supplied by the caller
(Listing~\ref{lst:market-listing}). Tying the badge to on-chain reputation, not to
a self-asserted parameter, is what stops a malicious source from minting fake
credibility.

\begin{lstlisting}[language=Solidity,caption={Verified flag sourced from the registry (\code{Marketplace.sol}).},label={lst:market-listing}]
listings[listingId] = Listing({
    whistleblower: msg.sender,
    descriptionHash: _descriptionHash,
    arweaveTxId: _arweaveTxId,
    minimumBid: _minimumBid,
    isActive: true,
    isVerified: registry.isVerified(msg.sender), // not self-asserted
    createdAt: block.timestamp,
    bidCount: 0
});
\end{lstlisting}

Settlement is where the anonymity guarantee is enforced. When a source accepts a
bid, they supply a freshly derived one-time stealth address; the contract deducts
a fixed platform fee and forwards the remainder there
(Listing~\ref{lst:market-accept}). Because the payout lands at an address with no
on-chain link to the source's published identity, an observer sees only a transfer
to an unfamiliar account.

\begin{lstlisting}[language=Solidity,caption={Paying out to a one-time stealth address under checks--effects--interactions (\code{Marketplace.sol}).},label={lst:market-accept}]
function acceptBid(uint256 _listingId, uint256 _bidIndex, address payable _stealthAddress)
    external nonReentrant
{
    Listing storage listing = listings[_listingId];
    require(listing.whistleblower == msg.sender, "Not listing owner");
    require(listing.isActive, "Listing not active");

    Bid storage bid = bids[_listingId][_bidIndex];
    require(!bid.isAccepted, "Bid already accepted");
    require(!bid.isWithdrawn, "Bid was withdrawn");
    require(_stealthAddress != address(0), "Invalid stealth address");

    // Effects before interactions.
    bid.isAccepted = true;
    listing.isActive = false;

    uint256 fee    = (bid.amount * PLATFORM_FEE_BPS) / 10000;
    uint256 payout = bid.amount - fee;

    (bool feeOk, )    = feeRecipient.call{ value: fee }("");
    require(feeOk, "Fee transfer failed");
    (bool payoutOk, ) = _stealthAddress.call{ value: payout }("");
    require(payoutOk, "Payout transfer failed");

    emit BidAccepted(_listingId, _bidIndex, _stealthAddress, payout);
}
\end{lstlisting}

The fund-moving functions --- \code{placeBid}, \code{acceptBid}, and
\code{withdrawBid} --- all follow the checks--effects--interactions pattern and are
additionally guarded with OpenZeppelin's \code{ReentrancyGuard} as
defence-in-depth. State is mutated to its final value \emph{before} any external
call, so a malicious recipient contract that re-enters during the value transfer
finds nothing left to exploit. This combination --- correct ordering plus a reentrancy
lock --- is the standard mitigation for the class of bug behind some of the largest
losses in the history of Ethereum, and it is applied uniformly here.

\section{Client-Side Cryptography}
\label{sec:impl-encryption}

All confidentiality in \zkw rests on a single primitive applied correctly:
AES-256-GCM, executed in the browser through the Web Crypto API. The
\code{encryption} service is the only place plaintext is ever handled, and it is
written so that plaintext never crosses a network boundary.

A fresh 256-bit key is generated per vault and marked \emph{extractable}, because
the key must later be exported so that Lit can seal it; it is decidedly not
persisted anywhere in the clear. Encryption draws a fresh 96-bit initialisation
vector --- the size recommended for GCM --- for every payload, guaranteeing that the
nonce is never reused under a given key (Listing~\ref{lst:encrypt}).

\begin{lstlisting}[language=TypeScript,caption={Per-payload key and IV, authenticated encryption (\code{encryption.ts}).},label={lst:encrypt}]
const ALGORITHM = "AES-GCM";
const KEY_LENGTH = 256;
const IV_LENGTH = 12; // 96-bit IV recommended for AES-GCM

export async function encryptData(data: Uint8Array, existingKey?: CryptoKey)
  : Promise<EncryptedPayload> {
  const key = existingKey ?? (await generateSymmetricKey());
  const iv  = crypto.getRandomValues(new Uint8Array(IV_LENGTH));

  const ciphertext = await crypto.subtle.encrypt({ name: ALGORITHM, iv }, key, data);
  const exportedKey = await exportKey(key); // raw bytes, to be sealed by Lit

  return { ciphertext: new Uint8Array(ciphertext), iv, exportedKey };
}
\end{lstlisting}

Because GCM is an authenticated mode, decryption fails loudly if the ciphertext or
IV has been tampered with; the system therefore gets integrity for free alongside
confidentiality, and a corrupted Arweave fetch cannot silently yield wrong
plaintext. The service also provides a compact binary serialisation,
\code{[4-byte IV length][IV][ciphertext]}, used when assembling the stored
manifest. Notably, the serialised form deliberately \emph{excludes} the key: the
exported key bytes travel only to Lit for sealing, never into storage.

This is the concrete realisation of principle P2 from
Chapter~\ref{ch:architecture}: the plaintext lives only in volatile browser memory
for the duration of encryption, and the key lives there only long enough to be
sealed. Everything that subsequently leaves the device is ciphertext.

\section{Threshold Key Custody with Lit Protocol}
\label{sec:impl-lit}

The most consequential implementation decision in \zkw is how the AES key is held
between arming and release. Chapter~\ref{ch:architecture} argued for replacing
hardware-rooted trust (a TEE) with mathematically-rooted trust (a threshold MPC
network). Lit Protocol provides exactly this: a decentralised network whose nodes
collectively hold key shares and will only reconstruct a secret when a stated
access-control condition evaluates true.

The condition \zkw uses is the dead man's switch itself: the helper that builds
it, \code{buildDeadMansSwitchACC}, produces an EVM-contract condition that reads
\isdeceased for the source's address and requires it to be \code{true}
(Listing~\ref{lst:lit-acc}).

\begin{lstlisting}[language=TypeScript,caption={Binding key release to the on-chain trigger (\code{litProtocol.ts}).},label={lst:lit-acc}]
export function buildDeadMansSwitchACC(contractAddress, userAddress, chain) {
  return [{
    conditionType: "evmContract",
    contractAddress,
    chain,
    functionName: "isDeceased",
    functionParams: [userAddress],
    functionAbi: {
      name: "isDeceased", type: "function", stateMutability: "view",
      inputs:  [{ name: "_user", type: "address", internalType: "address" }],
      outputs: [{ name: "",      type: "bool",    internalType: "bool" }],
    },
    returnValueTest: { key: "", comparator: "=", value: "true" },
  }];
}
\end{lstlisting}

Sealing the key is a client-side BLS encryption: anyone can encrypt a value to a
set of conditions without authenticating, because the conditions gate only
\emph{decryption}. Unsealing, by contrast, requires a wallet-authenticated session
(Sign-In With Ethereum) that the nodes verify, after which each node independently
checks the access-control condition against the chain and contributes a decryption
share only if it holds true. The combined shares reconstruct the key on the
requester's device; no single node, and no central server, ever sees the whole key.
This is the property a TEE cannot offer: trust is distributed across an
$n$-of-$m$ threshold rather than concentrated in one chip and its manufacturer.

Two implementation realities are worth recording, because they shaped the system.
First, Lit evaluates EVM conditions only against a fixed allow-list of public
chains; the local Hardhat chain (id \code{31337}) is not on it. The service
therefore exposes \code{isLitSupportedChain} and \code{litChainNameFromId}, and the
front end degrades gracefully to a ``local-preview'' mode on unsupported networks
so the rest of the flow remains demonstrable. Second, the project targets the Lit
SDK v8 stack (\code{@lit-protocol/lit-client} on the Naga network): the earlier v7
SDK and the Datil network it spoke to were sunset on 25~February~2026 and are no
longer reachable. Pinning to the MPC (Naga) path, rather than to newer
TEE-backed Lit runtimes, keeps the system aligned with the threshold-cryptography
trust model that motivates the whole design --- a point Chapter~\ref{ch:evaluation}
returns to.

\section{Permanent, Censorship-Resistant Storage}
\label{sec:impl-storage}

Encrypted vault manifests are stored on Arweave, whose one-time-fee,
pay-once-store-forever model is a better fit for a ``release on my death''
guarantee than a subscription-based pinning service that lapses if no one keeps
paying. Uploads go through a hardened Next.js server route rather than directly
from the browser, which lets the project use Turbo's free tier for sub-100~KiB
payloads --- no wallet, no funding transaction, no testnet gas --- while keeping the
upload path under the application's control.

This is the result of a deliberate design revision. An earlier iteration uploaded
directly from the browser using the Irys web SDK, which required the source's
connected wallet to pay for storage: the client had to query a price, check its
balance on the Irys bundler, and submit a \emph{funding transaction} before the
upload could proceed. For a whistleblowing tool this is doubly undesirable --- it
demands that the source hold and spend tokens, and, worse, it leaves an on-chain
funding trail that can become a deanonymisation vector. The system was therefore
migrated to ArDrive's Turbo SDK, whose \code{unauthenticated} free tier stores
payloads under 100~KiB with no wallet, no funding step, and no traceable
transaction. The migration also relocated the upload from the client to a
server-side route, which is what makes it possible to confine the Node-oriented
Turbo SDK to the server and to apply the hardening described below. The net effect
is a simpler, more robust path that also removes a privacy hazard --- a good example
of a security property emerging from an engineering simplification rather than from
added machinery. The module retains its original filename (\code{irysUpload.ts})
for continuity, but its implementation now wraps Turbo end to end.

The client side is deliberately thin: it POSTs raw bytes and a little metadata,
and reads payloads back by transaction id from a public gateway
(Listing~\ref{lst:storage-client}).

\begin{lstlisting}[language=TypeScript,caption={Thin storage client (\code{irysUpload.ts}).},label={lst:storage-client}]
export async function uploadToArweave(data: Uint8Array, metadata): Promise<string> {
  const response = await fetch("/api/storage", {
    method: "POST",
    headers: {
      "Content-Type": "application/octet-stream",
      "x-zkw-file-name": encodeURIComponent(metadata.fileName),
      "x-zkw-mime-type": metadata.mimeType,
      "x-zkw-version":   metadata.version,
    },
    body: data,
  });
  if (!response.ok) throw new Error(/* generic message */);
  const { id } = await response.json();
  return id;
}
\end{lstlisting}

The server route is where the security controls live (Listing~\ref{lst:storage-route}).
It runs only in the Node runtime --- so the Node-oriented Turbo SDK never executes
client-side --- caps the body at 100~KiB to enforce the free tier and act as an abuse
ceiling, applies best-effort per-IP rate limiting, attaches Arweave tags for
discoverability, returns generic error messages, and sets
\code{Cache-Control: no-store} on every response so that ciphertext is never
cached by intermediaries. Because the body is already AES-256-GCM ciphertext
wrapped in a self-describing manifest, the route handles no plaintext and no key
material; even a fully compromised storage server learns nothing.

\begin{lstlisting}[language=TypeScript,caption={Hardened upload route (abridged) (\code{api/storage/route.ts}).},label={lst:storage-route}]
export const runtime = "nodejs";
const MAX_BYTES = 100 * 1024;

export async function POST(request: Request) {
  const ip = request.headers.get("x-forwarded-for")?.split(",")[0]?.trim() || "unknown";
  if (rateLimited(ip)) return Response.json({ error: "Too many requests." },
                                            { status: 429, headers: noStore });

  const body = new Uint8Array(await request.arrayBuffer());
  if (body.byteLength === 0)          return Response.json({ error: "Empty payload." },        { status: 400, headers: noStore });
  if (body.byteLength > MAX_BYTES)    return Response.json({ error: "Payload too large." },    { status: 413, headers: noStore });

  const { TurboFactory } = await import("@ardrive/turbo-sdk");
  const turbo = TurboFactory.unauthenticated({ token: "base-usdc" });
  const { id } = await turbo.uploadRawX402Data({ data: Buffer.from(body), tags });
  return Response.json({ id }, { headers: noStore });
}
\end{lstlisting}

\section{Anonymous Payments: ERC-5564 Stealth Addresses}
\label{sec:impl-stealth}

To reward sources without doxxing them, \zkw implements ERC-5564 stealth addresses
over the secp256k1 curve using the audited \code{@noble/secp256k1} library. A
source publishes a \emph{stealth meta-address} --- a spending public key and a viewing
public key. A payer derives a fresh, one-time address from it for each payment;
only the source, using their private keys, can recognise and later spend the funds.

Derivation by the payer follows the standard construction: generate an ephemeral
key pair, compute a Diffie--Hellman shared secret against the source's viewing key,
hash it to a scalar, and add that scalar times the generator to the source's
spending public key (Listing~\ref{lst:stealth-derive}). The resulting point yields
the one-time Ethereum address; the first byte of the hashed secret is published as
a ``view tag'' so the source can skip non-matching announcements cheaply during
scanning.

\begin{lstlisting}[language=TypeScript,caption={Payer-side stealth-address derivation (\code{stealthAddress.ts}).},label={lst:stealth-derive}]
export function deriveStealthAddress(metaAddress): StealthPaymentInfo {
  const ephemeralPrivateKey = secp.utils.randomSecretKey();
  const ephemeralPublicKey  = secp.getPublicKey(ephemeralPrivateKey, true);

  // Shared secret: ephemeral_priv * viewing_pub
  const viewingPub   = hexToBytes(metaAddress.viewingPublicKey);
  const sharedSecret = secp.getSharedSecret(ephemeralPrivateKey, viewingPub, true);
  const hashedSecret = sha256(sharedSecret);

  // stealth_pub = spending_pub + hash(shared_secret) * G
  const spendingPub  = secp.Point.fromHex(metaAddress.spendingPublicKey);
  const offset       = secp.Point.BASE.multiply(bytesToBigInt(hashedSecret));
  const stealthPub   = spendingPub.add(offset).toBytes(false);

  return {
    stealthAddress: publicKeyToAddress(stealthPub),
    ephemeralPublicKey: bytesToHex(ephemeralPublicKey),
    viewTag: hashedSecret[0],
  };
}
\end{lstlisting}

Scanning mirrors this on the source's side: for each announced ephemeral key the
source computes the same shared secret with their viewing \emph{private} key, checks
the view tag, and confirms that the re-derived address matches the announced one.
The viewing key thus enables detection without exposing the spending key, which
remains the sole authority that can move the funds. The cryptographic symmetry
between the two paths --- payer multiplies by the viewing public key, source
multiplies by the viewing private key, and both arrive at the same secret --- is what
makes the link between payer and payee invisible on-chain while keeping the funds
fully spendable.

\section{Zero-Knowledge Provenance with Reclaim}
\label{sec:impl-reclaim}

Credibility and anonymity normally pull in opposite directions. \zkw reconciles
them with Reclaim Protocol, a zkTLS scheme that lets a source prove a fact derived
from an authenticated HTTPS session --- that they hold a particular corporate email
account, say --- without revealing the underlying credentials or their identity. The
proof's hash, or for the trustless path its on-chain-verified identifier, is
anchored in the registry; the bulky proof JSON stays off-chain.

The implementation's security hinge is that the Reclaim \emph{app secret} must
never reach the browser. \zkw enforces this with a server route that is the sole
holder of the secret: it performs \code{ReclaimProofRequest.init} server-side and
returns only a serialised, secret-free request configuration, which the client
rebuilds with \code{fromJsonString} (Listing~\ref{lst:reclaim-route}). The route
additionally validates the requested provider against a server-side allow-list,
rate-limits per IP, never logs or echoes secret material, and sets
\code{no-store}.

\begin{lstlisting}[language=TypeScript,caption={Server-held Reclaim secret (abridged) (\code{api/reclaim/route.ts}).},label={lst:reclaim-route}]
const ALLOWED_PROVIDER_IDS = new Set(Object.values(RECLAIM_PROVIDERS));

export async function POST(request: Request) {
  const appId     = process.env.RECLAIM_APP_ID;
  const appSecret = process.env.RECLAIM_APP_SECRET; // never sent to the client
  if (!appId || !appSecret)
    return Response.json({ error: "Verification is not configured." }, { status: 503, headers: noStore });

  const { providerId } = await request.json();
  if (typeof providerId !== "string" || !ALLOWED_PROVIDER_IDS.has(providerId))
    return Response.json({ error: "Unknown verification provider." }, { status: 400, headers: noStore });

  const reclaimProofRequest = await ReclaimProofRequest.init(appId, appSecret, providerId);
  return Response.json({ reclaimRequest: reclaimProofRequest.toJsonString() }, { headers: noStore });
}
\end{lstlisting}

On the client, the service starts the session, surfaces a request URL for the user
to complete the proof, and exposes a deterministic \code{hashProof} that produces a
canonical \code{keccak256} digest suitable for on-chain submission. For deployments
with a deployed verifier, the full proof can instead be transformed for on-chain
use and passed to \code{submitVerifiedProof} (Section~\ref{sec:impl-registry}),
giving the verification cryptographic, rather than merely attested, force.

\section{Front-End Orchestration}
\label{sec:impl-frontend}

The components and hooks of the \code{nextjs} package compose the services above
into two end-to-end flows: arming a vault and releasing it. The vault-creation
wizard is the clearest illustration of how the security properties chain together,
because it performs the pipeline in a fixed, auditable order: encrypt locally, then
seal the key with Lit, then upload only ciphertext, then register on-chain
(Listing~\ref{lst:vault-pipeline}).

\begin{lstlisting}[language=TypeScript,caption={The arming pipeline, in order (\code{VaultCreationForm.tsx}).},label={lst:vault-pipeline}]
// (1) Seal the AES key with Lit under the on-chain isDeceased() condition.
const acc    = buildDeadMansSwitchACC(deadMansSwitch.address, address, litChain);
const litKey = await encryptKey(encryptedPayload.exportedKey, acc, litChain);

// (2) Assemble a self-describing manifest -- never any plaintext.
const manifest: VaultManifest = {
  version: "1", app: "ZK-Whistle", encryption: "AES-256-GCM",
  file: { name: fileName, mimeType, size: encryptedPayload.ciphertext.length, encryptedAt: Date.now() },
  payload: { iv: bytesToBase64(encryptedPayload.iv),
             ciphertext: bytesToBase64(encryptedPayload.ciphertext) },
  lit: { network: LIT_NETWORK, chain: litChain,
         ciphertext: litKey.ciphertext, dataToEncryptHash: litKey.dataToEncryptHash,
         accessControlConditions: acc },
  recipient: recipientAddr,
};

// (3) Persist ciphertext to Arweave, then (4) register the switch on-chain.
const arweaveTxId = await upload(new TextEncoder().encode(JSON.stringify(manifest)), /* meta */);
await createSwitch(BigInt(interval), arweaveTxId, litRef, recipientAddr);
\end{lstlisting}

The ordering is the security argument made executable. Encryption precedes every
egress, so no plaintext can leak even if a later step fails; the key is sealed
before storage, so the manifest written to Arweave contains only ciphertext and a
\emph{sealed} key that is useless until the switch triggers; and registration comes
last, so a half-built vault is never advertised on-chain. The manifest is
self-describing, which means release needs nothing but the Arweave transaction id
recorded in the switch.

Release, implemented in the \code{useVaultRelease} hook, runs the pipeline in
reverse: fetch and parse the manifest from Arweave, authorise with the wallet and
ask Lit to release the AES key (which only succeeds once \isdeceased is true on a
supported chain), then AES-GCM-decrypt the payload locally and hand back the
recovered file. Decryption, like encryption, happens entirely on the requester's
device. It is worth being candid, as the code comments are, that this decrypt path
follows the documented v8 \code{AuthManager} pattern but has not yet been exercised
against a live Naga deployment; Chapter~\ref{ch:evaluation} treats this honestly as
a limitation.

\section{Deployment and Network Configuration}
\label{sec:impl-deploy}

Deployment uses hardhat-deploy scripts that run in dependency order: a mock Reclaim
verifier (for local testing), then the registry, then the marketplace, with the
dead man's switch independent of the others. The interesting logic lives in how
verification is wired per network (Listing~\ref{lst:deploy-registry}). On the local
chain the registry is deployed with \code{autoVerifyOnSubmit} enabled and pointed
at the mock verifier, so the full flow is demonstrable offline. On live networks
\code{autoVerifyOnSubmit} is disabled and the registry is pointed at the real,
per-chain Reclaim verifier address (overridable by environment variable), so
verification carries cryptographic weight.

\begin{lstlisting}[language=TypeScript,caption={Per-network verification wiring (abridged) (\code{01\_deploy\_whistleblower\_registry.ts}).},label={lst:deploy-registry}]
const isLocal = chainId === "31337";

await deploy("WhistleblowerRegistry", {
  from: deployer,
  args: [isLocal],   // autoVerifyOnSubmit: dev-only convenience
  log: true, autoMine: true,
});

let verifier = isLocal
  ? (await get("MockReclaim")).address
  : process.env.RECLAIM_VERIFIER_ADDRESS || RECLAIM_VERIFIERS[chainId];

if (verifier)
  await execute("WhistleblowerRegistry", { from: deployer }, "setReclaimVerifier", verifier);
\end{lstlisting}

After deployment, SE-2 regenerates the typed contract bindings consumed by the
front end, so the React hooks always reflect the deployed ABIs. The same scripts
deploy cleanly to a Lit-supported testnet such as Base Sepolia, which is the
configuration under which the threshold-custody path can be exercised end-to-end.

\section{Security Engineering Summary}
\label{sec:impl-security}

Security in \zkw is not a single feature but a set of invariants upheld at every
layer; Table~\ref{tab:impl-security} collects the principal controls and where they
are enforced. The throughline is least exposure: secrets are generated as late and
destroyed as early as possible, every network surface sees only ciphertext or
public data, and trust is distributed --- across the Lit threshold network, the
Reclaim witnesses, and the chain --- rather than concentrated in any one component.

\begin{table}[h]
\centering
\caption{Security controls and their enforcement points.}
\label{tab:impl-security}
\begin{tabular}{p{4.6cm} p{8.8cm}}
\toprule
\textbf{Control} & \textbf{Where and how enforced} \\
\midrule
In-browser encryption & AES-256-GCM via Web Crypto; plaintext never leaves the device (\code{encryption.ts}). \\
Fresh nonce per payload & 96-bit IV drawn per encryption; no key/nonce reuse (\code{encryption.ts}). \\
Threshold key custody & Key sealed under on-chain conditions; released only by Lit MPC quorum (\code{litProtocol.ts}). \\
No plaintext on-chain & Contracts store only timestamps, hashes, and references (\code{*.sol}). \\
Reentrancy safety & Checks--effects--interactions plus \code{ReentrancyGuard} on fund moves (\code{Marketplace.sol}). \\
Honest verification badge & Verified flag read from registry, not self-asserted; on-chain proof check (\code{Marketplace.sol}, \code{WhistleblowerRegistry.sol}). \\
Secret containment & Reclaim app secret confined to a Node-only route; provider allow-list (\code{api/reclaim}). \\
Storage hardening & Size cap, per-IP rate limit, generic errors, \code{no-store} (\code{api/storage}). \\
Payment unlinkability & One-time ERC-5564 stealth addresses per payout (\code{stealthAddress.ts}). \\
\bottomrule
\end{tabular}
\end{table}

With the system built, the next chapter turns to whether it works and how well:
the automated test suite, the deployment exercises, an evaluation against the
threat model and design goals, and a frank accounting of the prototype's
limitations.
